\documentclass{ctexart}
\usepackage{enumitem}
\usepackage{amsmath}
\newcommand{\zhquote}[1]{“#1”} % 定义了一个中文引号的命令
\begin{document}
\section{充分条件和必要条件}
命题是能够判断真假的陈述句:真的命题称为真命题,假的命题称为假命题.
下面的这些是真命题:
\begin{enumerate}
    \item $1>0$;
    \item 正方形是平行四边形；
    \item 平行直线间没有交点.
\end{enumerate}
下面是假命题的例子:
\begin{enumerate}
    \item $1 \le 0$;
    \item 方程 $x^2+1 =0$ 有实数解；
\end{enumerate}

上面介绍的都是最简单的命题.我们可以把简短的命题组合成一个长命题,呈“如果$\cdots $，那么$\cdots$”的形式.“$\cdots$”的位置可以用具体的命题代替,数学家们喜欢用$p$ 和 $q$表示命题，填入句子里就是“如果$p$，那么$q$”\footnote{数学家喜欢偷懒,数学就是一门用精简的方式描述我们所在的世界的学科}.

让我们来看具体的例子:
\zhquote{如果$x$是整数,那么$x$是有理数};这里$p$命题就是\zhquote{$x$是整数},$q$命题是\zhquote{$x$是有理数}.

可以发现\zhquote{如果$x$是整数,那么$x$是有理数}这句话是对的,所以这个命题是真命题.
若 “如果$p$，那么 $q$ ”是真命题,即 $p$ 可以由 $q$ 推导出,则称 $p$ 是 $q$ 的充分条件,记作$ p \Rightarrow q$.反之如果是假命题,即$p$ 不能推导出 $q$ ,则称 $p$ 不是 $q$ 的充分条件,记作$p \not\Rightarrow q $.

必要条件是和充分条件对偶的概念, $ p \Rightarrow q$时, $p$ 是 $q$ 的充分条件,而$q$ 是 $p$ 的必要条件.\textbf{口诀}:箭头出发的是另一个充分条件,箭头所指的是另一个的必要条件.
\[\text{\textbf{充分条件}}\Longrightarrow  \text{\textbf{必要条件}}\]



\section{平面向量}
\subsection{线性运算}



\end{document}